\section{Выбор паттернов проектирования}
\headingtextsep

При проектировании системы выбраны три паттерна: Repository, Adapter и Strategy. В контексте реализации на Rust UML-диаграммы классов используются как средство описания структур, трейтов и сервисных компонентов, выполняющих роли классов и интерфейсов.

\textheadingsep
\subsection{Паттерн Repository}
\headingtextsep

Паттерн Repository применяется для изоляции прикладной логики от деталей хранения данных. Бизнес-модули работают с абстракциями репозиториев, а конкретные реализации инкапсулируют обращение к PostgreSQL и поисковым индексам.

Примером применения является работа с задачами, документами и заметками, где прикладные сервисы используют интерфейсы репозиториев вместо прямых SQL-запросов. UML-диаграмма класса для этого паттерна приведена в приложении А.

\textheadingsep
\subsection{Паттерн Adapter}
\headingtextsep

Паттерн Adapter используется для подключения внешних инфраструктурных зависимостей через внутренние абстракции системы. Это позволяет не связывать доменную модель с конкретным API провайдера и упрощает замену внешнего компонента без переработки прикладной логики.

Основной пример применения в системе --- абстракция AssetStorage, скрывающая различия между локальным файловым хранением и S3-compatible backend. UML-диаграмма класса для этого паттерна приведена в приложении А.

\textheadingsep
\subsection{Паттерн Strategy}
\headingtextsep

Паттерн Strategy применяется там, где система должна поддерживать несколько взаимозаменяемых вариантов поведения. В архитектуре системы такой подход целесообразен для выбора механизма аутентификации человека в зависимости от режима работы и подключенного провайдера.

Примером применения является выбор между GitHub OAuth, OIDC и dev-auth режимом через единый интерфейс провайдера аутентификации. UML-диаграмма класса для этого паттерна приведена в приложении А.